%\section{General Notions on Continuous Groups}
%\subsection{Continuous Groups and Infinitesimal Groups}

We start with a set of \(n\) variables \(x_0^i\), \(i=1 \ldots n\), which may be
regarded as coordinates of a point in a certain space. Consider now the set of equations
\begin{equation}\label{eq:1}
x^i=f^i\left(x_0^1, \ldots, x_0^n ; a^1, \ldots, a^r\right) \quad(i=1 \ldots n),
\end{equation}
in which the \(a^\rho\) appear as a set of \(r\) independent parameters.
Omitting indices, we shall write this and similar relations in the form
\[
x=f\left(x_0 ; a\right) \quad \text { or } \quad x=S_a x_0 \tag{1'}
\]

These equations define a set \(S\) of transformations, depending on the
parameters \(a\), which map the point \(x_0\) onto \(x\). We shall assume
that the \(f^i\) have all the required derivatives, and that the \(f^i\) depend
\underline{essentially}
on the parameters, i.e., that no two transformations with different
parameters are the same for all values of \(x_0\), so that \(r\) is the
smallest number of parameters needed to specify the transformations completely
and uniquely.


The set of transformations \(f\) is said to form a group if it obeys the
following two conditions:
\begin{itemize}
\item [i)]
The result of performing successively any two transformations of the set is
another transformation belonging to this set. Formally, if
\(x=f\left(x_0 ; a\right)\) and \(x^{\prime}=f(x ; b)\)
then there exists a set of parameters \(c^\rho\)
\begin{equation}\label{eq:2}
c^\rho=\varphi^\rho(a ; b)
\end{equation}
such that
\begin{equation}\label{eq:3}
x^{\prime}=f(x ; b)=f\left(f(x_0; a) ; b\right)=f(x_0 ; c)=f(x_0 ; \varphi(a ; b)).
\end{equation}
\item [ii)]
Corresponding to every transformation there exists a unique inverse, which also
belongs to the set: Given equation (\ref{eq:1}) there exists a set of
parameters \(\bar{a}\) such that \(x_0=f(x ;\bar{a})\).
\end{itemize}

The uniqueness of \(\bar{a}\) is guaranteed if the Jacobian of the transformation does not vanish:
\begin{equation}\label{eq:4}
\left|\frac{\partial f}{\partial x_0}\right| \neq 0
\end{equation}

Transforming \(x_0\) onto \(x\) and then inversely \(x\) back to \(x_0\), we
obtain according to i) a transformation which belongs to the group and is
characterized by the set of parameters \(a_0\). Since the transformation
depends on the parameters in an essential way, the \(a\) so constructed cannot
depend on the particular value of the parameters from which we started. The
transformation \(f(x, a_0)\) is called the identity.

Since it imposes no restriction, we shall take:
\[
a_0^\rho=\textcolor{red}{0} \quad (\rho=1 \ldots r). \quad 
\]
\(r\) is called the \underline{order} of the group. (Note that this usage is
different from that found in the theory of finite groups.)

We also remind the reader of the following definitions: A mapping of one group
onto another is said to be \underline{homomorphic} or a
\underline{homomorphism} if it preserves the operation of group multiplication.
We call such a mapping an \underline{isomorphism} if, in addition, the
correspondence between elements of the two groups is one-to-one.  Since the
combination law of the transformations (\ref{eq:1}) is given in terms of the
parameters \(a\), there can be transformations corresponding to different
values of \(n\) which are homomorphic or even isomorphic.

A group of linear transformations which is homomorphic with a given group is called a
\underline{representation} of this group.

The fundamental idea of Sophus Lie's theory of continuous groups~\cite{Lie1893} is to consider
not the whole of a group, but that part of it which lies near the identity,
consisting of the so-called \underline{infinitesimal}
\underline{transformations}. Thus, instead of the finite displacement of a
point under a transformation, we consider the application of successive
infinitesimal displacements - we think of a generalized velocity field
describing the motion of a point from its original position \(x_0\) to its
final position \(x\).
\textcolor{blue}{The velocity field is independent of the parameters.}


We have now two equivalent expressions for x:
\begin{subequations}\label{eq:5}
  \begin{align}
x=f(x_0; a), \label{eq:5a} \\
x=f(x; 0).  \label{eq:5b}
  \end{align}
\end{subequations}
Corresponding to these we can represent in either of two ways a transformation,
as a result of which the new components of \(x\) differ infinitesimally from
the old ones - by differentiation of (\ref{eq:5a}) or by introducing a parameter
of infinitesimal size in (5b):
\[
x+d x=f\left(x_0 ; a+d a\right) \quad \text{or} \quad
x+d x=f\left(x; \delta a\right)
\]
or (employing the summation convention),
\begin{equation}\label{eq:6}
d x=\frac{\partial f\left(x_{0}; a\right)}{\partial a^\sigma} d{a}^\sigma
\quad \text{or} \quad
d x=\left(\frac{\partial f(x ; a)}{\partial a^\sigma}\right)_{a=0} \delta a^\sigma.
\end{equation}
The last may be written
\begin{equation}\label{eq:7}
d x^i=u_\sigma^i(x) \delta a^\sigma, \quad
u_\sigma^i(x)=\left(\frac{\partial f^i(x; a)}{\partial a^\sigma}\right)_{a=0}
\end{equation}
which defines the ``velocity field'' \(u_\sigma^i(x)\) mentioned above.
\textcolor{blue}{The velocity field is independent of the parameters.}
In the notation of (\ref{eq:2}) we may write
\[
a+d a=\varphi(a ; \delta a) .
\]
Since it follows from (\ref{eq:2}) and (\ref{eq:5}) that \(\varphi(a ; 0)=a\), we have
\[
a+d a=a+\left(\frac{\partial \varphi(a ; b)}{\partial b^\tau}\right)_{b=0} \delta a^\tau .
\]
Thus, da is a linear combination of \(\delta a\) :
\begin{equation}\label{eq:8}
d a^\rho=\mu_\tau^\rho(a) \delta a^\tau, \quad \mu_\tau^\rho(a)=
\left(\frac{\partial \varphi^\rho(a; b)}{\partial b^{\tau}}\right)_{b=0}
\end{equation}

Solving for \(\delta a\), we get
\[
\delta a^\sigma=\lambda_{\rho}^{\sigma} (a) ~~ d {a}^{\rho} \tag{8'}
\]
where
\begin{equation}\label{eq:9}
\lambda \mu=1, \qquad \text{i.e.}~~
\lambda_{\rho}^\sigma \mu_\tau^\rho=\delta_\tau^\sigma.
\end{equation}
From (6), (7) and (8') we get the first fundamental formula:
\[
\frac{\partial x^i}{\partial a^{\rho}} = u^i_{\tau}(x) \lambda^{\tau}_{\rho} (a) \tag{A}
\]
If \(u\) is to represent the velocity field of a transformation (\ref{eq:1}), equation
(A) must be completely integrable, i.e. it must be capable of admitting
solutions with \(n\) arbitrary constants \(x_0\). The integrability condition
\[
\frac{\partial^2 x^i}{\partial a^{\sigma} \partial a^{\rho}}=
\frac{\partial^2 x^i}{\partial a^{\rho} \partial a^{\sigma}}
\]
\textcolor{blue}{
\[
\frac{\partial^2 x^i}{\partial a^{\sigma} \partial a^{\rho}}=
\left( \frac{\partial}{\partial a^{\sigma}} u^i_{\nu} \right) \lambda^{\nu}_{\rho}+
u^i_{\nu} \frac{\partial \lambda^{\nu}_{\rho}}{\partial a^{\sigma}} 
\]
\[
\frac{\partial u^i_{\nu} }{\partial a^{\sigma}}  =
\frac{\partial u^i_{\nu} }{\partial x^{j}} 
\frac{\partial x^{j}}{\partial a^{\sigma}}, \quad
\frac{\partial x^j}{\partial a^{\sigma}} = 
u^j_{\tau} \lambda^{\tau}_{\sigma}
\]
\[
\frac{\partial u^i_{\nu} }{\partial a^{\sigma}} \lambda^{\nu}_{\rho}=
\frac{\partial u^i_{\nu} }{\partial x^{j}} 
u^j_{\tau} \lambda^{\tau}_{\sigma}
 \lambda^{\nu}_{\rho}
\]
}
becomes
\textcolor{red}{
\[
\left( u_{\kappa}^j \frac{\partial u^i_{\nu}}{\partial x^j}-
u^j_{\nu} \frac{\partial u_{\kappa}^i}{\partial x^j}\right)
\lambda_{\sigma}^{\kappa} \lambda^{\nu}_{\rho} +
u^i_{\nu} \left( \frac{\partial \lambda_\sigma^\nu}{\partial a^{\rho}}-
\frac{\partial \lambda_{\rho}^{\nu}}{\partial a^{\sigma}} \right)=0.
\]
}
and, using (\ref{eq:9}),
\textcolor{blue}{
\[
d_{\kappa \nu}^{i} := 
u_{\kappa}^j \frac{\partial u^i_{\nu}}{\partial x^j}-
u^j_{\nu} \frac{\partial u_{\kappa}^i}{\partial x^j}, \quad
e^{\tau}_{\rho  \sigma} := 
\frac{\partial \lambda_{\rho}^{\nu}}{\partial a^{\sigma}} -
 \frac{\partial \lambda_\sigma^\nu}{\partial a^{\rho}}
\]
\[
\lambda^{\nu}_{\rho} \mu^{\rho}_{q} =\delta^{\nu}_{q}, \quad
\lambda^{\kappa}_{\sigma} \mu^{\sigma}_{m} =\delta^{\kappa}_{m}, \quad
d_{\kappa \nu}^{i} \lambda_{\sigma}^{\kappa}
\lambda^{\nu}_{\rho} \mu^{\sigma}_{m} \mu^{\rho}_{q} = d_{q m}^{i} =
u^i_{\tau} e^{\tau}_{\rho  \sigma} \mu^{\sigma}_{m} \mu^{\rho}_{q} 
\]
}
this gives \textcolor{blue}{equation B1,}
\begin{equation}\label{eq:10}
u_{\kappa}^j \frac{\partial u^i_{\nu}}{\partial x^j}-
u^j_{\nu} \frac{\partial u_{\kappa}^i}{\partial x^j} =
c_{\kappa \nu}^{\tau}(a)  u_{\tau}^{i}
\end{equation}
where
\begin{equation}\label{eq:11}
c_{\kappa \nu}^{\tau}(a) = \left(
\frac{\partial \lambda^{\tau}_{\rho}}{\partial a^{^\sigma}}-
\frac{\partial \lambda_{\sigma}^{\tau}}{\partial a^\rho}
\right)
\mu_{\kappa}^\rho \mu_{\nu}^{\sigma}.
\end{equation}
Since \(u\) is independent of \(a\), differentiation of (10) by \(a^{\rho}\) gives
\[
\frac{\partial c_{x \nu}^\tau(a)}{\partial a^{\rho}} u_\tau^i=0 .
\]
But the a's have been assumed essential, so that by (7) the u's are
linearly independent, hence the c's are independent of \(a\).  From (\ref{eq:11})
\[
\frac{\partial \lambda_{\rho}^{\tau}}{\partial a^{\sigma}} -
 \frac{\partial \lambda_\sigma^\tau}{\partial a^{\rho}} =
c_{\kappa \nu}^{\tau} \lambda_\rho^{\kappa} \lambda_\sigma^\nu. \tag{B2}
\]
(B1) is a necessary condition on the velocity field if the latter is to
generate a group, and (B2) is a corresponding restriction on the manner in which
the a's combine.

An infinitesimal transformation on the \(x\) induces on any function \(F(x)\) a variation
\begin{equation}\label{eq:12}
d F(x)=\frac{\partial F}{\partial x^i} d x^i=
\delta a^\sigma u_\sigma^i \frac{\partial F}{\partial x^i} \equiv \delta a^\sigma x_\sigma F
\end{equation}
where
\begin{equation}\label{eq:13}
x_\sigma=u_\sigma^i(x) \frac{\partial}{\partial x^i}
\end{equation}
(\ref{eq:12}) shows that every infinitesimal transformation of \(F(x)\) is generated by
a linear combination of the operators \(X\) which are called the infinitesimal
operators of the group \(S\). From \(\left(B_1\right)\) it follows that they
satisfy the relation
\begin{equation}\label{eq:14}
x_\rho x_\sigma-x_\sigma x_\rho \equiv\left[x_\rho x_\sigma\right]=c_{\rho \sigma}^\tau x_\tau .
\end{equation}
Evidently
\[
c_{\rho \sigma}{ }^\tau=-c_{\sigma \rho}{ }^\tau .
\tag{C1}
\]
Substituting (14) into the Jacobi identity
\[
\left[\left[X_\rho X_\sigma\right] X_\tau\right]+
\left[\left[X_\sigma X_\tau\right] X_\rho\right]+
\left[\left[X_\tau X_\rho\right] X_\sigma\right]=0
\]
we get
\[
c_{p \sigma}{ }^\mu c_{\mu \tau}{ }^\nu+c_{\sigma \tau}^\mu c_{\mu p}^\nu+c_{\tau \rho}{ }^\mu c_{\mu \sigma}{ }^\nu=0.
\tag{C2}
\]
We have shown that equations (c) are implied if the \(f^i\) form a group. That the converse to this statement holds is the content of the three fundamental theorems of Lie~\cite{Lie1893}, which we shall not prove. They state that
\begin{itemize}
\item[I.] If there exist \(f^i=x^i\) satisfying (A) then they form a group.
\item[II.] If there exist u's satisfying \(\left(B_1\right)\), then there exist \(\lambda\) 's, determined within isomorphism, which satisfy \(\left(B_2\right)\), so that equation (A) is integrable.
\item[III.] For every set of o's satisfying (C), there exist u's satisfying \(\left(B_1\right)\).
\end{itemize}
We shall write an infinitesimal transformation of the group \(S\) in the form
\[
S_a=1+\delta a^\sigma X_\sigma,
\]
where \(\delta a^\sigma\) is an infinitesimal quantity defined to be of
the first order. If we combine two such transformations, we get
\[
S_a S_b= \left(1+\delta a^{\rho} X_\rho\right)\left(1+\delta a^{\sigma} X_{\sigma}\right)=
1+\delta a^{\rho} X_{\rho}+\delta a^{\sigma} X_{\sigma},
\]
where the first non-vanishing infinitesimal terms have been retained.
Thus, to the operation of multiplication in \(S\) corresponds addition in the
infinitesimal group of \(S\). If the first order quantities vanish, we have to
consider quantities of higher order. But the second theorem of Lie implies that
in this connection we need never go beyond the second order of infinitesimals -
i.e. we have only to worry about commutators, which are expressions of the form
\(S_a S_b S_a^{-1} S_b^{-1}\), and to ask that the corresponding infinitesimal
operator of the second order,
\[\delta a^{\rho} \delta b^\sigma\left[X_\rho X_\sigma\right]
\]
be contained in the linear manifold of infinitesimal operators.
